 %! TeX program = lualatex
\documentclass[10pt,a4paper]{article}

	\usepackage{amsmath,amssymb,amsthm,mathtools,amsfonts}
	\usepackage{breqn}
	\usepackage[utf8]{inputenc}
	\usepackage[danish]{babel}
	\usepackage[T1]{fontenc}

	%%%%%%%%%%%%%%%%%%%%%%%%%%%%%%%%%%%%%%%%%%%%%%%%%%%%%%%%%%%%%%%%%%%%%%%%%%%%%% Fonts

	\usepackage{fontspec}
	\setmainfont[Numbers=OldStyle]{EB Garamond}
	\newfontfamily\plantinc[Numbers=OldStyle]{Plantin MT Pro Bold Condensed}
	\newfontfamily\plantinsb{Plantin MT Pro Semibold}
	\newfontfamily\garamond[Numbers=OldStyle]{EB Garamond}
	\newfontfamily\plantin[Numbers=OldStyle]{Plantin MT Pro}
	\newfontfamily\wal{Walbaum Com}
	\newfontfamily\klav{Klavika}
	\newfontfamily\klavl{Klavika light}
	\newfontfamily\quotefont[Ligatures=TeX]{Linux Libertine O}

	%%%%%%%%%%%%%%%%%%%%%%%%%%%%%%%%%%%%%%%%%%%%%%%%%%%%%%%%%%%%%%%%%%%%%%%%%%%%%% Colors

	\usepackage{xcolor}
  \usepackage{transparent}

	\definecolor{frbblue}{RGB}{47, 88, 109}
	\definecolor{hgred}{RGB}{140, 10, 10}
	\definecolor{frbgreenl}{RGB}{162, 202, 137}
	\definecolor{frbgreen}{RGB}{28, 85, 66}
	\definecolor{frbgrey}{RGB}{243, 245, 242}


	%%%%%%%%%%%%%%%%%%%%%%%%%%%%%%%%%%%%%%%%%%%%%%%%%%%%%%%%%%%%%%%%%%%%%%%%%%%%%% Other Graphics

	\usepackage{graphicx}
	\graphicspath{{./img/}}
	\usepackage{tikz}
	\usepackage{placeins}
	\usepackage[modulo]{lineno}

	\usepackage[pdfborder={0 0 0}]{hyperref} %% Ingen firkant rundt om referencer

	\usepackage{multicol}
	\usepackage{cellspace}
		\setlength\cellspacetoplimit{4pt}
		\setlength\cellspacebottomlimit{4pt}
	\usepackage{tabularx}
		\newcolumntype{b}{>{\hsize=\hsize}>{\centering\arraybackslash}X}
	\usepackage{siunitx}

	\usepackage{enumitem}
	\usepackage{wrapfig}
		\setlength{\columnsep}{0pt}
		\setlength{\intextsep}{-10pt}
	\usepackage{caption}

	\setlength{\parskip}{0.5\baselineskip}
	\setlength{\parindent}{0cm}

	\setlist[enumerate]{label=\alph*),left=20pt}

	\usepackage{lipsum}
	\usepackage{comment}

	%%%%%%%%%%%%%%%%%%%%%%%%%%%%%%%%%%%%%%%%%%%%%%%%%%%%%%%%%%%%%%%%%%%%%%%%%%%%%% New commands

	\newcommand\svar[1]{\newline\textcolor{red}{\textbf{Svar}\\#1}}

	\usepackage[h]{esvect}
	\newcommand{\twvect}[2]{\ensuremath{\begin{pmatrix}#1\\#2\end{pmatrix}}}

	\newcommand*\quotesize{20}
	\newcommand*{\openquote}
		{\tikz[remember picture,overlay,xshift=-3ex,yshift=-0ex]
			\node (OQ) {\quotefont\fontsize{\quotesize}{\quotesize}\selectfont``};\kern0pt}

	\newcommand*{\closequote}[1]
		{\tikz[remember picture,overlay,xshift=2ex,yshift={#1}]
			\node (CQ) {\quotefont\fontsize{\quotesize}{\quotesize}\selectfont''};}

	\newcommand{\qm}[1]{``#1"}
	\newcommand{\iquote}[1]{\begin{quote}\itshape \openquote#1\hfill\closequote{0ex} \end{quote}}

	\newcommand{\source}[1]{\footnote{Source: \href{#1}{#1}}}

	\newcommand{\glos}[3]{\dotuline{#1}\marginpar{\scriptsize\textit{#3} #2}}

	\newcommand{\subtitle}[1]{\vspace{10pt}\newline\normalsize\plantin #1}

	%%%%%%%%%%%%%%%%%%%%%%%%%%%%%%%%%%%%%%%%%%%%%%%%%%%%%%%%%%%%%%%%%%%%%%%%%%%%%% Sections

	\usepackage{titlesec}
	\titleformat{\subsection}{\color{hgred}\plantinc\large}{}{}{}

	\usepackage{titling}
	\setlength{\droptitle}{-12ex}
	\pretitle{\begin{flushleft}\plantinc\huge\color{frbblue}}
	\posttitle{\par\end{flushleft}}
	\preauthor{\begin{flushleft}\Large}
	\postauthor{\end{flushleft}}
	\predate{\begin{flushleft}}
	\postdate{\end{flushleft}}

	\setcounter{secnumdepth}{0}

	%%%%%%%%%%%%%%%%%%%%%%%%%%%%%%%%%%%%%%%%%%%%%%%%%%%%%%%%%%%%%%%%%%%%%%%%%%%%%% Header & Footer

	\usepackage{bigfoot}
	\DeclareNewFootnote{a}
		\renewcommand{\thefootnotea}{\fnsymbol{footnotea}}
			\newcommand{\footnotes}[1]{\footnotea[2]{#1}}
			\newcommand{\footnoted}[1]{\footnotea[3]{#1}}
			\MakePerPage{footnotes}

			%%%%%%%%%%%%%%%%%%%%%%%%%%%%%%%%%%%%%%%%%%%%%%%%%%%%%%%%%%%%%%%%%%%%%%%%%%%%%% Document

\begin{document}
%\pagecolor{FloralWhite}
\title{My Name}
%\date{\vspace{-3em}} % I tilfældet ingen dato
\date{1984}
\author{Sandra Cisneros}

\maketitle

\noindent\linenumbers In English my name means hope. In Spanish it means too many letters. It means sadness, it means waiting. It is like the number nine. A muddy color. It is the Mexican records my father plays on Sunday mornings when he is shaving, songs like sobbing.

It was my great-grandmother's name and now it is mine. She was a horse woman too, born like me in the Chinese year of the horse — which is supposed to be bad luck if you're born female — but I think this is a Chinese lie because the Chinese, like the Mexicans, don't like their women strong.

My great-grandmother. I would've liked to have known her, a wild horse of a woman, so wild she wouldn't marry. Until my great-grandfather threw a sack over her head and carried her off. Just like that, as if she were a fancy chandelier. That's the way he did it.

And the story goes she never forgave him. She looked out the window her whole life, the way so many women sit their sadness on an elbow. I wonder if she made the best with what she got or was she sorry because she couldn't be all the things she wanted to be. Esperanza. I have inherited her name, but I don't want to inherit her place by the window.

At school they say my name funny as if the syllables were made out of tin and hurt the roof of your mouth. But in Spanish my name is made out of a softer something, like silver, not quite as thick as sister's name --- Magdalena --- which is uglier than mine. Magdalena who at least can come home and become Nenny. But I am always Esperanza.

I would like to baptize myself under a new name, a name more like the real me, the one nobody sees. Esperanza as Lisandra or Maritza or Zeze the X. Yes. Something like Zeze the X will do.

\bigskip
\noindent\textit{Source: Sandra Cisneros, \textit{The House on Mango Street} (1984), Vintage Books.}

\end{document}
